% vim: set ft=tex tabstop=4 shiftwidth=4 noexpandtab:

\input{../../include/latex/tikz/header.tex}

% opening %{{{1

\begin{document}
\begin{tikzpicture}[scale=1.0]

% parameters %{{{1

	\tikzmath{
		\radius = 3.5;
	}

% coordinates %{{{1

	\coordinate (O_1) at (0,0);
	\coordinate (O_2) at (\radius,0);

% intersections %{{{1

	\begin{pgfinterruptboundingbox}
		\path[name path=circleA] (O_1) circle (\radius);
		\path[name path=circleB] (O_2) circle (\radius);
		\path[name intersections={of=circleA and circleB, by={I,J}}];
	\end{pgfinterruptboundingbox}

	\begin{pgfinterruptboundingbox}
		\path[name path=line1] (O_1) -- (O_2);
		\path[name path=line2] (I) -- (J);
		\path[name intersections={of=line1 and line2, by=M}];
	\end{pgfinterruptboundingbox}

% distance %{{{1

	%\path let \p1 = ($ (B) - (A) $)
	%in \pgfextra{
		%\pgfmathsetmacro{\distAB}{veclen(\x1,\y1)/1cm}
		%\global\let\distAB\distAB
	%};

% points, dots, vertices %{{{1

	\fill (O_1) circle (0.4mm);
	\fill (O_2) circle (0.4mm);
	\fill (I) circle (0.4mm);
	\fill (J) circle (0.4mm);
	\fill (M) circle (0.4mm);

% segments, sides, lines %{{{1

	\draw (O_1) -- (O_2);
	\draw (I) -- (J);
	\draw (O_1) -- (I);
	\draw (O_1) -- (J);
	\draw (O_2) -- (I);
	\draw (O_2) -- (J);

% circles %{{{1

	\draw (O_1) circle (\radius);
	\draw (O_2) circle (\radius);

% circular arcs %{{{2

	%\draw (A) ++(\ang:\radius) arc (\ang:\gle:\radius);

% points, dots, vertices labels %{{{1

	\node[left] at (O_1) {$O_1$};
	\node[right] at (O_2) {$O_2$};
	\node[above] at (I) {$I$};
	\node[below] at (J) {$J$};
	\node[above left] at (M) {$M$};

% segments, sides, lines labels %{{{1

	\node[above left] at ($ (O_1)!0.5!(I) $) {$r$};
	\node[below left] at ($ (O_1)!0.5!(J) $) {$r$};

	\node[above right] at ($ (O_2)!0.5!(I) $) {$r$};
	\node[below right] at ($ (O_2)!0.5!(J) $) {$r$};

	\node[below] at ($ (O_1)!0.5!(M) $) {$a$};
	\node[below] at ($ (O_2)!0.5!(M) $) {$b$};

	\node[right] at ($ (M)!0.4!(I) $) {$c$};
	\node[right] at ($ (M)!0.4!(J) $) {$d$};

% segments, sides, lines marks %{{{1

	\tkzMarkSegments[mark=|, size=2pt, pos=0.5](O_1,M)
	\tkzMarkSegments[mark=|, size=2pt, pos=0.5](O_2,M)

	\tkzMarkSegments[mark=||, size=2pt, pos=0.5](O_1,I)
	\tkzMarkSegments[mark=||, size=2pt, pos=0.5](O_1,J)

	\tkzMarkSegments[mark=||, size=2pt, pos=0.5](O_2,I)
	\tkzMarkSegments[mark=||, size=2pt, pos=0.5](O_2,J)

	\tkzMarkSegments[mark=|, size=2pt, pos=0.35](M,I)
	\tkzMarkSegments[mark=|, size=2pt, pos=0.40](M,I)
	\tkzMarkSegments[mark=|, size=2pt, pos=0.45](M,I)

	\tkzMarkSegments[mark=|, size=2pt, pos=0.35](M,J)
	\tkzMarkSegments[mark=|, size=2pt, pos=0.40](M,J)
	\tkzMarkSegments[mark=|, size=2pt, pos=0.45](M,J)

% circles labels %{{{1

	\node at ($ (O_1) + (145:\radius + 0.3) $) {$\mathscr{C}_1$};
	\node at ($ (O_2) + (35:\radius + 0.3) $) {$\mathscr{C}_2$};

% circular arcs labels %{{{2

	%\node at ($ (O) + (45:\radius + 0.3) $) {$\mathscr{C}$};

% angles labels %{{{1

	%\pic[draw, ->, "$\alpha$", angle radius=0.8cm, angle eccentricity=1.5]
		%{angle = B--A--C};

% right angles markers %{{{2

	%\pic [draw, angle radius=7] {right angle = O_2--M--I};

% closing %{{{1

\end{tikzpicture}
\end{document}
